\section{Анализ путей в графе и линейная алгебра}
\label{sec:PathAlgebra}
\tikzsetfigurename{GraphTheory_PathAlgebra_}

Часто поиск путей сопровождается изучением их свойств, что далее приводит к формулированию дополнительных ограничений на пути в терминах этих свойств.
Например, можно потребовать, чтобы пути были простыми или не проходили через определённые вершины.
Один из естественных способов описывать свойства и, как следствие, задавать ограничения~--- это использовать ту алгебраическую структуру, из которой берутся веса рёбер графа%
\sidenote{На самом деле здесь наблюдается некоторая двойственность.
    С одной стороны, действительно, удобно считать, что свойства описываются в терминах некоторой заданной алгебраической структуры.
    Но, вместе с этим, структура подбирается исходя из решаемой задачи.}.

Предположим, что дан граф $\mbfscrG = \langle V, E, L\rangle $, где над $L$ построено полукольцо $\algebraic{L} = (S, \opAdd, \opMult)$\label{itm:otimesIntro}.
Тогда изучение свойств путей можно описать следующим образом:
\begin{equation}
    \label{eq:algPathProblem}
    \left\{\ (v_i, v_j, c) \mid c = \bigoplus_{v_i \pi_k v_j} \bigotimes_{(u, l, v) \in \pi } l \ \right\}.
\end{equation}

Иными словами, для каждой пары вершин, для которой существует хотя бы один путь, их соединяющий, мы агрегируем (с помощью операции $\opAdd$ из полукольца) информацию обо всех путях между этими вершинами.
При этом информация о пути получается как свёртка меток рёбер пути с использованием операции%
\sidenote{Заметим, что детали свёртки вдоль пути зависят от свойств полукольца (и от решаемой задачи).
    Так, если полукольцо коммутативно, то нам не обязательно соблюдать порядок рёбер.
    В дальнейшем мы увидим, что данные особенности полукольца существенно влияют на особенности алгоритмов решения соответствующих задач.}
$\opMult$.

Естественным требованием (хотя бы для прикладных задач, решаемых таким способом) является существование и конечность указанной суммы.
На данном этапе мы не будем касаться того, какие именно свойства полукольца могут нам обеспечить данное свойство, однако в дальнейшем будем считать, что оно выполняется.
Более того, будем стараться приводить частные для конкретной задачи рассуждения, показывающие, почему это свойство выполняется в рассматриваемых в задаче ограничениях.

Описанная выше задача общего вида называется анализом свойств путей алгебраическими методами (Algebraic Path Problem)~\sidecite{Baras2010PathPI} и предоставляет общий способ для решения широкого класса прикладных задач%
\sidenote{В работе \enquote{Path Problems in Networks}~\cite{Baras2010PathPI} собран действительно большой список прикладных задач с описанием соответствующих полуколец.
    Сводная таблица на страницах 58--59 содержит 29 различных прикладных задач и соответствующих полуколец.}.
Наиболее известными являются такие задачи, как построение транзитивного замыкания графа и поиск кратчайших путей (All Pairs Shortest Path или APSP).
Далее мы подробнее обсудим эти две задачи и предложим алгоритмы их решения.


 В данной главе мы рассмотрим некоторые связи%
 \sidenote{Связь между графами и линейной алгеброй~--- обширная область, в которой можно даже выделить отдельные направления, такие как спектральная теория графов.
     С точки зрения практики данная связь также подмечена давно и более полно с ней можно ознакомиться, например, в работах~\cite{doi:10.1137/1.9780898719918, Davis2018Algorithm9S}.
     Кроме этого, много полезной информации можно найти на сайте \url{https://graphblas.github.io/GraphBLAS-Pointers/}.}
 между графами и операциями над ними и матрицами и операциями над матрицами.

 Как мы видели, достаточно естественное представление графа~--- это его матрица смежности.
 Далее можно заметить некоторое сходство между определением матричного умножения~\ref{def:MxM} и мыслями, которыми мы руководствовались, вводя операцию $\opMult$~\ref{itm:otimesIntro}.
 Действительно, пусть есть матрица $M$ размера $n \times n$ над $\algebraic{G} = (S, \opAdd, \opMult)$~--- матрица смежности графа.
 Умножим её саму на себя: вычислим $M'= M \mmultFrom{G} M$.
 Тогда, по~\ref{def:MxM}, $M'[i, j] = \bigoplus_{l \in [0 \rng n-1]} M[i, l] \opMult M[l, j]$.
 Размер $M'$ также $n \times n$.
 То есть $M'$ задаёт такой граф, что в нём будет путь со свойством, являющимся агрегацией свойств всех путей, составленных из двух подпутей в исходном графе.
 А именно, если в исходном графе есть путь из $i$ в $l$ с некоторым свойством (его значение хранится в $M[i, l]$), и был путь из $l$ в $j$ (его значение хранится в $M[l,j]$), то в новом графе свойство $M[i, l] \opMult M[l, j]$ аддитивно (используя $\opAdd$) учтётся в свойстве пути из $i$ в $j$.

 \begin{example}
    \label{exmpl:matrixMultAndPaths}
     \NiceMatrixOptions{code-for-first-row = \color{red},
                     code-for-first-col = \color{blue},
                     code-for-last-row = \color{red},
                     code-for-last-col = \color{blue}
                     }

    Проиллюстрируем связь матричного умножения и вычисления свойств путей.
    Пусть матрица $A$ содержит информацию о рёбрах из стартовых вершин в промежуточные, а $B$~--- из промежуточных в конечные.
    Тогда $C = A \cdot B$ агрегирует информацию о путях длины~$2$, составленных из двух подпутей.
     Схематично, вычисление элемента $c_{ij}$ показано ниже ($c_{ij} = \bigoplus_{k} a_{ik} \opMult b_{kj}$):
\[
\begin{pNiceArray}{cccc}[first-col,last-col,nullify-dots]
  &  & \Cdots &  &  & \\
i & a_{i0} & \Cdots & & a_{i(k-1)} & i \\
  &  & \Cdots &  &  & \\
  &  & \Cdots &  &  &
\end{pNiceArray}
\times
\begin{pNiceArray}{cccc}[first-row,last-row=5,nullify-dots]
 & & j &                             \\
   &              & b_{0j} &         \\
  \Vdots & \Vdots & \Vdots & \Vdots  \\
   &              &        &         \\
   &              & b_{(k-1)j} &     \\
 & & j &
\end{pNiceArray}
=
\begin{pNiceArray}{cccc}[first-col,last-col,first-row,last-row=5,nullify-dots]
   &         &        & j       &        &   \\
   &         &         & \Vdots  &        &   \\
 i & \Cdots  & \Cdots  & c_{ij} & \Cdots & i \\
   &          &        & \Vdots &        &  \\
   &          &        & \Vdots &        &  \\
    &         & & j       &        &
\end{pNiceArray}
\]

    Соответствующий фрагмент графа для двух промежуточных вершин $k_0$ и $k_1$ изображён на рисунке~\ref{fig:algebraic_path_problem_path_example}.

 \begin{marginfigure}
    \begin{tikzpicture}[shorten >=1pt,auto]
    \node[state,fill=blue!20] (q_0)                      {$i$};
    \node[state] (q_1) [right = of q_0 ]    {$k_0$};

    \node[state] (q_2)  [below = of q_1 ]    {$k_1$};
    \node[state,fill=red!20] (q_3)  [right = of q_1]    {$j$};

    \path[->]
    (q_0) edge  node {$a_0$} (q_1)
    (q_0) edge[right]  node {$a_1$} (q_2)
    (q_1) edge  node {$b_0$} (q_3)
    (q_2) edge[right]  node {$b_1$} (q_3);


 \path[->,red]
   (q_0) edge[bend left=40]  node {$ a_0 \otimes b_0 \oplus a_1 \otimes b_1 $} (q_3)
   ;

\end{tikzpicture}

\caption{Вычисление свойств пути при умножении матриц смежности}
\label{fig:algebraic_path_problem_path_example}
\end{marginfigure}
 \end{example}

 Таким образом, произведение матриц смежности соответствует обработке информации о путях интересующим нас образом.
 Это наблюдение позволяет предложить решение задач анализа свойств путей, основанное на операциях над матрицами.
 Проиллюстрируем этот подход на двух классических задачах: транзитивном замыкании и поиске кратчайших путей.

 Рассмотрим сперва более простую задачу~--- транзитивное замыкание графа, которая также является частным случаем~\eqref{eq:algPathProblem}.
 Здесь метки рёбер не несут дополнительной информации: нас интересует лишь факт достижимости.
 Соответствующее полукольцо~--- булево: $\algebraic{B} = \langle \{0,1\}, \lor, \land \rangle$, где $\Bbbzero = 0$ (пути нет), $\Bbbone = 1$ (путь существует).
 Матрица смежности $M$ над $\algebraic{B}$ определяет отношение смежности вершин: $M[i,j] = 1$ тогда и только тогда, когда существует ребро из $i$ в $j$.

 Матричное умножение над $\algebraic{B}$ принимает вид
 \[
     (M \cdot N)[i,j] = \bigvee_{k \in [0 \rng n-1]} M[i,k] \land N[k,j].
 \]
 Иными словами, $(M \cdot N)[i,j] = 1$ в точности тогда, когда существует вершина~$k$, для которой $M[i,k] = N[k,j] = 1$.

 Следовательно, $M^d$ содержит информацию о путях длины ровно~$d$: $M^d[i,j] = 1$ тогда и только тогда, когда существует путь из $i$ в $j$, состоящий ровно из $d$ рёбер.


 \begin{example}
     \label{exmpl:transitiveClosure}
     Рассмотрим граф из примера~\ref{exmpl:diGraph}.

     \begin{marginfigure}
         \resizebox{\marginparwidth}{!}{\begin{tikzpicture}[]
  % Apex angle 60 degrees
  \node[isosceles triangle,
      isosceles triangle apex angle=60,
      draw=none,fill=none,
      minimum size=2cm] (T60) at (3,0){};

  \node[state] (q_0) at (T60.right corner) {$0$};
  \node[state] (q_1) at (T60.left corner)  {$1$};
  \node[state] (q_2) at (T60.apex)         {$2$};
  \node[state] (q_3) [right=1.5cm of q_2]  {$3$};
  \path[->]
    (q_0) edge (q_1)
    (q_1) edge (q_2)
    (q_2) edge (q_0)
    (q_2) edge [bend left] (q_3)
    (q_3) edge [bend left] (q_2);
  \path[->, red]
    (q_0) edge [bend left] (q_2)
    (q_1) edge [bend right] (q_0)
    (q_1) edge [bend left] (q_3)
    (q_2) edge [bend right] (q_1)
    (q_2) edge [loop above] (q_2)
    (q_3) edge [bend left] (q_0)
    (q_3) edge [loop right] (q_3);
\end{tikzpicture}
}
         \caption{Исходные рёбра графа (чёрным) и пути длины~$2$ (красным)}
         \label{fig:transitiveClosure_paths_length2}
     \end{marginfigure}

     Его булева матрица смежности~--- $M$, а квадрат этой матрицы~--- $M^2$:
     \[
         M = \begin{pmatrix}
             0 & 1 & 0 & 0 \\
             0 & 0 & 1 & 0 \\
             1 & 0 & 0 & 1 \\
             0 & 0 & 1 & 0
         \end{pmatrix},\qquad
         M^2 = \begin{pmatrix}
             0 & 0 & 1 & 0 \\
             1 & 0 & 0 & 1 \\
             0 & 1 & 1 & 0 \\
             1 & 0 & 0 & 1
         \end{pmatrix}.
     \]
     Результат этого шага представлен на рисунке~\ref{fig:transitiveClosure_paths_length2}, где красным цветом показаны рёбра, соответствующие путям длины~$2$ в исходном графе.
     Видим, что, например, $M^2[0,2] = 1$, так как существует путь $0 \to 1 \to 2$ длины~$2$.
     При этом $M^2[0,0] = 0$, поскольку путь длины ровно~$2$ из $0$ в $0$ отсутствует, а рефлексивности нет (в исходной матрице на диагонали стоят нули).
 \end{example}

 В задаче транзитивного замыкания нас интересуют пути любой длины, а не только фиксированной.
 Для этого необходимо замкнуть отношение смежности относительно транзитивности.
 Добавим петли во все вершины с весом $\Bbbone$: положим $M' = I \lor M$, где $I$~--- единичная матрица ($I[i,i] = \Bbbone$, $I[i,j] = \Bbbzero$ при $i \neq j$).
 Тогда $(M')^d$ содержит информацию о путях длины \emph{не более}~$d$~\sidenote{Работает линейная алгебра: $(M+I)^2 = M^2 + IM + MI + I^2$. То есть появляются пути длины два и остаются длины один. Далее можно продолжать в том же духе.}.
 Поскольку в ориентированном графе ни один простой путь не содержит более $n-1$ ребра, транзитивное замыкание соответствует $(M')^{n-1}$, причём суммировать промежуточные степени не нужно~--- они уже содержатся в старшей.

 Таким образом, задача транзитивного замыкания свелась к возведению булевой матрицы в степень%
 \sidenote{Рекомендуем построить транзитивное замыкание для примера~\ref{exmpl:transitiveClosure}.}
  $n-1$, что является частным случаем общего подхода~\eqref{eq:algPathProblem}.

 Заметим, что описанная схема~--- возведение матрицы смежности с петлями нейтрального веса в степень $n-1$ над подходящим полукольцом~--- является общей для широкого класса задач анализа свойств путей.
 Рефлексивность гарантирует, что $M^k$ содержит информацию о путях длины не более $k$, а ограниченность простых путей $n-1$ рёбрами~--- что степень $n-1$ достаточна.

 Применим эту схему к задаче APSP.
 Здесь используется тропическое полукольцо $\algebraic{T} = \langle \BbbR \cup \{+\infty\}, \min, + \rangle$, а рефлексивность обеспечивается нулями на главной диагонали матрицы смежности $M$.
 Тогда $M_k = M^k$~--- матрица кратчайших путей длины не более $k$, и, как и в случае транзитивного замыкания, достаточно вычислить $M_{n-1}$.
 Здесь мы также пользуемся тем, что в ориентированном графе без отрицательных циклов кратчайший путь не может содержать более $n$ рёбер%
 \sidenote{Раз отрицательных циклов нет, то проходить через одну вершину, коих $n$, больше одного раза бессмысленно.}.

 Таким образом, решение APSP сведено к произведению матриц над $\algebraic{T}$, однако вычислительная сложность наивного подхода слишком большая: $O(n \cdot \SMM{\algebraic{T}}{n})$.

 Благодаря рефлексивности%
 \sidenote{Как мы видели на примере транзитивного замыкания, матрица $M_k$ содержит информацию о путях длины не более $k$, а не ровно $k$.}%
 , для вычисления $M_{n-1}$ можно применить \emph{быстрое возведение в степень}.
 Будем последовательно возводить матрицу в квадрат, получая матрицы $M_2$, $M_4$, $M_8$ и так далее.

 В итоге получим следующую последовательность вычислений:
 \begin{gather*}
     M_1 = M \\
     M_2 = M_1 \cdot M_1 = M^2 \\
     M_4 = M_2 \cdot M_2 = M^4 \\
     M_8 = M_4 \cdot M_4 = M^8 \\
     \vdots \\
     M_{2^k} = M_{2^{k-1}} \cdot M_{2^{k-1}} = M^{2^k}
 \end{gather*}

 Здесь мы предполагаем, что $n-1 = 2^k$ для некоторого $k$.
 На практике такое верно далеко не всегда.
 Если это условие не выполняется, то необходимо взять ближайшую сверху степень двойки%
 \sidenote{
     Кажется, что это приведёт к избыточным вычислениям.
     Попробуйте оценить, на сколько много лишних вычислений будет сделано в худшем случае.}%
 , вычислить $M_{2^{\lceil \log(n-1) \rceil}}$, после чего результат будет содержать и искомую матрицу $M_{n-1}$.

 Теперь вместо $n$ итераций нам нужно $\log{n}$, а итоговая сложность решения~--- $O(\log{n} \cdot \SMM{\algebraic{T}}{n})$%
 \sidenote{
     Заметим, что это оценка для худшего случая.
     На практике при использовании данного подхода можно прекращать вычисления как только при двух последовательных шагах получились одинаковые матрицы ($M_i = M_{i-1}$).
     Это приводит нас к понятию \emph{неподвижной точки}, обсуждение которого лежит за рамками повествования.}%
 .
 Данный алгоритм%
 \sidenote{Пример решения APSP с помощью быстрого возведения в степень: \url{http://users.cecs.anu.edu.au/~Alistair.Rendell/Teaching/apac_comp3600/module4/all_pairs_shortest_paths.xhtml}.}
 называется \emph{быстрым возведением в степень} (repeated squaring).

 Итак, быстрое возведение в степень сводит APSP к $O(\log n)$ умножениям матриц над тропическим полукольцом.
 Однако, как обсуждалось в разделе~\ref{sec:TheoreticalComplexity}, быстрые алгоритмы матричного умножения (алгоритм Штрассена и его развитие, дающие оценку $O(n^\omega)$ с $\omega < 2.373$) требуют возможности вычитания, то есть работают над кольцами, но не над полукольцами.
 Для умножения матриц над полукольцами лучшие известные алгоритмы являются \emph{слегка субкубическими} (mildly subcubic): наилучшая оценка составляет $\hat{O}(n^3 / \log^4 n)$~\sidecite{10.1007/978-3-662-47672-7_89}.
 Таким образом, и APSP через этот подход получает лишь слегка субкубическую сложность $O(n^3 \log n / \log^4 n)$.

 Вопрос о существовании \emph{истинно} субкубического алгоритма для общего случая APSP остаётся открытым~\sidecite{Chan2010} и является одной из центральных проблем мелкозернистой сложности.
 Вместе с тем активно развиваются подходы, дающие истинно субкубические алгоритмы для структурированных частных случаев Min-Plus произведения и соответствующих подклассов APSP: Вильямс и Сюй (SODA'20)~\sidecite{3381089.3381091} предложили алгоритм для матриц, блоки которых близки к матрицам константного ранга; Сюй (2021)~\sidecite{xu_2021}~--- для матриц с константным аддитивным приближённым рангом и монотонных матриц; Чи, Дуан и Се (SODA'22)~\sidecite{doi:10.1137/1.9781611977073.60}~--- для ограниченно-разностных матриц с оценкой $O(n^{2.779})$.

 Заметим, что быстрое возведение в степень, использованное нами для транзитивного замыкания и APSP, возможно только при ассоциативности произведения матриц\sidenote{Ассоциативность гарантируется свойствами полукольца.}.
 Для более бедных алгебраических структур, не обеспечивающих ассоциативность, данный приём неприменим.

\mycomment{Резкий переход к разделу про GraphBLAS и программные абстракции. Нужен связующий абзац-мостик.}
